\section{Citation and license}
\label{sec:citation}

\subsection{Citing \denis{}}

\denis{} is developed in the open at
\href{https://github.com/ardakayaalp/DENIS}{github.com/ardakayaalp/DENIS}
and every release is archived on Zenodo. The version DOI for
release~1.0.0 is
\href{https://doi.org/10.5281/zenodo.22081267}{10.5281/zenodo.22081267};
the concept DOI
\href{https://doi.org/10.5281/zenodo.22081266}{10.5281/zenodo.22081266}
always resolves to the latest version. The canonical citation entry,
including machine-readable metadata, lives in the \code{CITATION.cff}
file at the repository root and is surfaced as a ``Cite this
repository'' button on the GitHub project page.

\paragraph{Suggested citation.}
\begin{quote}
    A.~Kayaalp, \emph{DENIS: Doppler Estimation and Numerical
    Inference for Spectroscopy}, version~1.0.0, Zenodo (2026).
    DOI:~\href{https://doi.org/10.5281/zenodo.22081267}%
    {10.5281/zenodo.22081267}.
\end{quote}

\paragraph{Acknowledged dependencies.}
\denis{} relies on satlas2~\cite{satlas2} for hyperfine-structure
fitting and on \code{clstools}~\cite{clstools} for IGISOL-style
ASDF I/O. Please cite these packages alongside \denis{} when
reporting analyses produced with the toolkit. For background on the
underlying technique, the Handbook of Nuclear Physics chapter on
nuclear spins and moments~\cite{cls_handbook} is a standard
reference.

\subsection{License}

\denis{} is released under the \textbf{MIT License} --- see the
\code{LICENSE} file at the repository root. You may use, modify, and
redistribute it freely, provided the copyright and permission notice
are retained. The bundled Lucide icons are MIT/ISC licensed;
\code{satlas2} and \code{clstools} carry their own licenses.

\subsection{Contact}

Bug reports, feature requests, and contributions are welcome via the
GitHub issue tracker. For correspondence about the underlying
physics or for collaboration enquiries, contact the developer at
\href{mailto:arda.kayaalp@kuleuven.be}{arda.kayaalp@kuleuven.be}.
